\documentclass[17pt, a4paper]{extarticle}
\usepackage{graphicx}
\usepackage{xcolor}
\usepackage{comment}
\usepackage{cite}
\usepackage{url}
\usepackage{hyperref}
\usepackage{wrapfig}
\usepackage[a4paper, total={6in, 8in}]{geometry}
\geometry{legalpaper, landscape, margin=2in}
\begin{document}

\begin{titlepage}
	\centering
	\includegraphics[width=0.1\textwidth]{logo}\par\vspace{1.0cm}
	{\scshape\LARGE IIT Bombay \par}
	\vspace{1.0cm}
	\color{green}
	{\LARGE\bfseries Dirt Bike Simulation\par}
	\color{black}
	\vspace{1.0cm}
	{\itshape Tejesh Raut\par}
	{\itshape Shrey Kumar\par}
	{\itshape Gangesh Gudmalwar\par}
	\vfill
	supervised by
	Prof.~Sharat \textsc{Chandran}

	\vfill
\end{titlepage}

\color{red}
\textbf{Honour Code}\newline
\newline
\color{black}
We pledge on our honour that we have not given or received any unauthorized assistance on this assignment or any previous task.\newline \newline
\color{red} 
\textbf{Project Idea}
\color{black}
\newline \newline
This is a Box2D simulation that has  a dirt bike which passes through various obstacles moving various things while itself moving.
The bike itself has movable parts which provides the shock-up features so that the bike can move through any rough track.
Basically, we designed a basic simulation of the using some features  of the popular dirt-bike game. It has a bike which is capable of going through any rough surfaces. It has a good grip and shock-up system which allows it to perform this task. This is the main part of the dirt-bike game and can be further used to make the actual dirt-bike game.
\newpage
\color{red} 
\textbf{Techniques used}\newline
\newline
\color{black}
At first we made the ground of our simulation by copying the code of ground of dominos simulation and changed the length of the ground so as to play with the bike for a long time.
\newline
Then we added the tilted ground by studying how the parameters are given while making the line and we got many tilted paths in the project.
\newline
Then we added many obstacles to the path such as pendulum, see-saw, garbage boxes and the three big boxes.
\newline
For the pendulum and see-saw we studied the dominos project and then copied the codes for these two objects.
\newline
For the garbage boxes we made them using loops and declared them dynamic and allowed them to fall on the tilted path from above. We made a small barrier in the tilted path to prevent them from sliding down from the slope.
\newline
We added the three big boxes one above the other at last to make a tuff game by crossing these boxes.
\newline
We then created the curved path by applying mathematical equations using loops
\newline
Then we made the bike using many dynamic objects linked to each other by joints.
We first made the whole body of the bike as a single part but it did not provide flexibility to the bike. So afterwards we made the bike by joining various parts of the bike by using joints. It allowed the bike to turn at the points where it had a joint.
\newline
Then we made the big heavy ball which initialized the motion of the bike by striking on its rear part
\newline
At last we decided to add a driver to the bike with its head turning front and back with the motion of the bike being motivated by the game of Hill-climb racing. For this we added a thin rectangular joined to the center part of the bike and then we added a small circular body to denote the head of the driver at the other end of the rectangular body.
\newline
For the outlab.pdf part we used github for maintaining the edits done in the project. We used the screenrecorder in ubuntu for recording the demo video. We tried hard but we were not able to record audio of the demonstration.
\newpage
\color{red} 
\textbf{Some screen shots}
\color{black}
\newline \newline
\begin{wrapfigure}{l}{0.5\textwidth}
\centering
	\includegraphics[width=0.5\textwidth]{fig1.png}
\end{wrapfigure}
\begin{wrapfigure}{l}{0.2\textwidth}
\centering
	\includegraphics[width=0.2\textwidth]{fig2.png}
\end{wrapfigure}
\color{red}
\textbf{Who has done what?}\newline
\newline
\color{black}
We coordinated among each other and everyone has participated in coding and design part. The distribution given below is just a rough idea of who participated on that part most.
\newline
Tejesh and Shrey did the code of the path and all small objects in the path. Curved path was created by Shrey. Shrey created the bike with Gangesh. We all created the driver of the bike. We all participated in resizing and repositioning of the objects in proper ratio and place. Gangesh and Tejesh assigned keyboard controls to the bike and changed the panning controls to the \textbf {a, s, d} and \textbf {w} which were originally controlled by the arrow keys. Arrow keys are assigned to control the direction and tilt of the bike.
\newline
Regarding the outlab.pdf, the part regarding maintaining the git repository was done jointly by Tejesh. Initially we did not maintain the edits in git but we started adding edits to git as soon as we came to know about maintaining a git repository. Task of editing the makefile was done jointly by Shrey and Gangesh. The project webpage was made by Tejesh and is available at here(Shrey yaha pe mere Box2D page ki link daal dena) . Code documentation was done by Gangesh. Profiling was done by Gangesh. This report was made jointly by Tejesh and Shrey.
\newpage
\color{red} 
\textbf{Deviations}\newline
\newline
\color{black}
There aren't any great deviations in the design of the project. But we added some more features to the project by adding driver to the bike and keyboard controls which we did not mention in our ProjectProposal.
\newline
We decided to include the keyboard controls to make it interesting so that a user can himself play the game instead of just seeing the simulation.
\newline
We were not able to add graphics to our project because as faar as we searched we need to add some packages in the Box2D project to allow graphics in the project.
\newpage
\color{red} 
\textbf{SRS}\newline
\newline
\newline
\color{blue}
\textbf{Purpose}\newline
\color{black}
The purpose of this project is to build a basic version of dirt-bike game. Many additions can be done to this project to make the actual dirt-bike game
\newline
\newline
\color{blue}
\textbf{Scope}\newline
\color{black}
For running this project you need to have linux in your PC and you must install the Box2D, the steps for installing of which are given in the .txt files in the src folder. This dirt-bike simulation game is not intended to work in extreme cases. It is just the basic version of the real dirt-bike game.
\newline
\newline
\color{blue}
\textbf{How the project works?}\newline
\color{black}
The simulation is initialized by dropping of the ball from a height at the top left corner. The ball strikes the bike from its backside giving a momentum to the bike. The bike then falls on the see-saw thus making the box present on the left side of the see-saw to bounce off. Then the bike goes through the garbage boxes after passing through the pendulum in the path. Finally the bike hits the three large boxes kept at the end over each other. After this you may cross the three big boxes present at last by playing tactfully smile emoticon .
\newline
\newline
\color{blue}
\textbf{Motivation}\newline
\color{black}
Our motivation was the inspiration of Box2D Mountain Bike as given in the outlab.pdf of lab3. We googled through various small Box2D games. We decided to make a simulation deriving from the game of dirt bike. We were also motivated by the game Hill-climb racing which is somewhat similar to the dirt-bike game.
\newline
\newline
\color{blue}
\textbf{Conclusion}\newline
\color{black}
This part of the project can be further used to make the complete game of dirtbike by changing the path as per the game. We can edit the features of tilting the bike leftwards or rightwards so that they may work better than present ones to get the feel of the real dirt-bike game.
\color{black}
\newpage
\color{red} 
\textbf{Observations regarding the profiling }\newline
\color{black}
As main part in project was bike, it has these many joints joining two wheels, driver and body part , so the SolveVelocityConstraints() takes more  a time 
and for these bike part we used relative co-ordinates to assign position so function operators involving vectors also takes more time  compared to
 original box2D prject .\newpage
\color{red} 
\textbf{Refrences}\newline
\color{black}
\cite{bib1}
\cite{bib2}
\bibliography{Report}
\bibliographystyle{plain}
\end{document}